\documentclass{article} % For LaTeX2e
\usepackage{iclr2026_conference,times}

% Optional math commands from https://github.com/goodfeli/dlbook_notation.
\input{math_commands.tex}

\usepackage{hyperref}
\usepackage{cleveref}
\usepackage{url}

\usepackage{amsmath}
\usepackage{amssymb}
\usepackage{booktabs}
\usepackage{graphicx}
\usepackage{subcaption}
\usepackage{xcolor}
\usepackage{xspace}
\usepackage{algorithm}
\usepackage{algpseudocode}

\newcommand{\note}[1]{\textcolor{red}{#1}}
\makeatletter
\newenvironment{ai}
  {\color{lightgray}%
   \@ifpackageloaded{caption}{\captionsetup{font={color=lightgray}}}{}%
   \let\ai@savedfloatreset\@floatboxreset
   \g@addto@macro\@floatboxreset{\color{lightgray}}}
  {\global\let\@floatboxreset\ai@savedfloatreset}
\makeatother

\input{macros.tex}

\graphicspath{{fig/}{fig/slides/}}

\title{Elucidating Inference in \\ Multi-Token Prediction}

% Authors must not appear in the submitted version. They should be hidden
% as long as the \iclrfinalcopy macro remains commented out below.
% Non-anonymous submissions will be rejected without review.
%
% TODO(draft): confirm the final author list and affiliations before submission.
\author{Justus Will\thanks{Correspondence: \texttt{jcwill@uci.edu}.} \\
Department of Computer Science\\
University of California, Irvine\\
Irvine, CA 92697, USA
}

\newcommand{\fix}{\marginpar{FIX}}
\newcommand{\new}{\marginpar{NEW}}

%\iclrfinalcopy % Uncomment for camera-ready version, but NOT for submission.
\begin{document}

\maketitle

\begin{abstract}
% Intro setting + emphasize impact
Accelerating generation is a key challenge in deployment of modern LLMs.
% Intro problem
Multi-token prediction provides substantial inference speedups by emitting several tokens per forward pass, but the quality of the resulting text depends critically on how those tokens are accepted, rejected, or resampled.
% Intro approach
We study this decoding layer directly, treating the multi-token model as a black-box proposal mechanism and identify design choices that best trade throughput against generation quality or alignment to an underlying autoregressive (AR) reference.
% Main point / insight
Together, our improvements yield up to 10 tokens per call across a wide range of applications on several pretrained models, a 200\% increase over recent methods, unlocking the full low-latency potential of inference on these models.
\end{abstract}

\begin{ai}
\section{Introduction}

\begin{itemize}
    \item AR LLM are the base of modern inference
    \item many methods to speed up including small draft models. MTP/PTP improves up to 3x.
    \item design choices have been focused on exactly replicating teacher, if not biased trade-off (example p-threshold)
    \item proposal of parallizaton (apple, PTP)
\end{itemize}

Autoregressive decoding emits one token per forward pass. The autoregressive
model is therefore both the bottleneck and the reference: it sets the latency of
a deployed language model, and it defines what the output is supposed to be,
since any method that accelerates it is judged by whether it still produces what
that model would have produced.

A long line of work accelerates that reference without replacing it. Speculative
decoding drafts with a small autoregressive model and verifies the draft against
the target~\citep{leviathan2023speculative,chen2023accelerating}; later methods
make the drafter cheaper by attaching extra decoding heads to the target
itself~\citep{stern2018blockwise,cai2024medusa,ankner2024hydra,li2024eagle}, or
drop the drafter entirely and retrieve candidate continuations from the
prompt~\citep{saxena2023pld}. Multi-token prediction pushes hardest on the
proposal, training a single model to emit several coordinated tokens per
call~\citep{gloeckle2024multitoken,draxler2025ptp}. Task-average speedups near
$2.4\times$, approaching $3\times$ on the most favourable task categories, are
now routine on diverse-task benchmarks~\citep{xia2024specbench}.

A multi-token model breaks the one-token-per-pass constraint by proposing several
tokens at once, but a proposal is not an output: something must decide which of
the proposed tokens are kept. That decision --- the \emph{acceptance rule} --- is
usually treated as a fixed implementation detail, inherited from speculative
decoding, where it is chosen to leave the base model's output distribution
exactly intact.

That exactness is already being traded away. Medusa~\citep{cai2024medusa} argues
that matching the base model's distribution is typically unnecessary and instead
accepts any token the base model finds sufficiently plausible, a rule
Hydra~\citep{ankner2024hydra} inherits. The trade is real and it is large: on our
checkpoint, switching the predicate from exact matching to nucleus membership at
$\theta = 0.98$ and changing nothing else raises committed tokens per call from
$3.63$ to $4.30$ and wall-clock speedup from $1.02\times$ to $1.70\times$, while
the mean probability the base model assigns to the generated tokens falls by
$17$ points. Once the proposal mechanism is strong enough to offer many coherent
tokens per call, the acceptance rule stops being an implementation detail: it
sets the throughput, and it sets how far the generated text drifts from what the
base model would have produced. Yet it is chosen by convention and reported by
nobody --- the standard protocol measures accepted tokens per call and speedup,
never faithfulness.

This paper takes the decoding layer as the object of study. We ask the question
a practitioner actually faces: \emph{on a single GPU, what is the lowest latency
we can achieve that still produces high-quality outputs?} We organize the answer
along three axes.

\textbf{Verification.} What counts as a correct token? The strictest rule accepts
a proposed token only if the base model, driven by the same randomness, would have
emitted exactly that token. Relaxing it to ``the base model finds this token
plausible'' --- nucleus membership, a probability threshold, a likelihood ratio
--- accepts far more tokens per call, at a cost in faithfulness that we measure
rather than assume.

\textbf{Proposal geometry.} How many chances to take, and how long each should
be. Given a node budget, is it better to propose one long continuation, or fifty
short ones? We answer this by modelling the number of correct tokens per call as
the extinction time of a branching process, which makes the expected yield
separable across depth and the allocation problem exactly solvable.

\textbf{Parallelization.} Verification and the next proposal can run
concurrently if one speculates on the verification outcome, which recent work
has begun to exploit~\citep{samragh2025future,draxler2025ptp}. Partial Quadratic
Decoding (PQD)~\citep{draxler2025ptp} does this by preparing a branch for each
possible number of correct tokens and spending a budget across those branches. We
show that the branch weights and the value of a branch can be read off the same
calibrated model that solves the geometry problem, rather than estimated
separately.

We make three main contributions.

\begin{itemize}
  \item \textbf{A design space that generalizes existing methods.} We separate
  multi-token decoding into three independent axes --- verification, proposal
  geometry, and parallelization --- and resolve the published methods along them
  (Table~\ref{tab:baselines}). Most are special cases: a chain is a width profile
  of one, a static draft tree is a profile fixed offline, and EAGLE-2's
  context-aware tree is a greedy solution to our allocation problem under an
  independence assumption that we make explicit and then remove. Eleven
  verification rules sit behind a single code path, so two configurations differ
  only in the acceptance predicate, and we pair them with a measurement protocol
  that separates throughput from faithfulness --- including a test of whether the
  generated text is distributionally distinguishable from the base model's
  (\S\ref{sec:design}).

  \item \textbf{A generative model of $\ncorrect$ that makes proposal geometry
  solvable.} A zero-inflated Beta-Binomial branching process with shared
  per-generation frailty predicts per-call accepted-token histograms across three
  orders of magnitude of proposal width from three parameters
  ($R^2 \geq 0.976$). Its survival function makes expected yield separable across
  depth, which turns strand-length selection into a knapsack we solve exactly by
  dynamic programming --- and the same object supplies the branch weights and
  branch values that Partial Quadratic Decoding needs, replacing two separately
  calibrated estimates with one (\S\ref{sec:model}--\S\ref{sec:pqd}).

  \item \textbf{Evidence that the standard metric is the wrong one alone.}
  Across $62$ decoding configurations on Spec-Bench, accepted-tokens-per-call,
  wall-clock speedup, and faithfulness to the base model induce three different
  orderings. The configurations at the top of the tokens-per-call ranking get
  there by accepting almost everything, roughly halving the base model's mean
  token probability; among those that stay close to the base model, more tokens
  per call often means \emph{slower} generation --- wide ensembles reach $4.6$
  tokens per call at $0.31\times$ autoregressive speed. Only the exact rules are
  distributionally indistinguishable from the base model
  (\S\ref{sec:experiments}).
\end{itemize}

\section{Background: multi-token proposals and their verification}
\label{sec:background}

\subsection{Auxiliary variables}

We build on Parallel Token Prediction (PTP)~\citep{draxler2025ptp}, which we
recap only as far as needed. A standard decoder predicts
$P_i := P(t_i \mid t_{<i})$ and then samples $t_i \sim P_i$. The sampling step is
where the sequential dependency lives, and it can be made explicit: drawing
$\aux_i \sim \mathcal{U}[0,1]$ and taking the inverse CDF,
\begin{equation}
  t_i = \Pick(\aux_i, P_i) \equiv \min\Big\{ j : \textstyle\sum_{l \leq j} P_{il} > \aux_i \Big\},
  \label{eq:pick}
\end{equation}
makes $t_i$ a \emph{deterministic} function of the context and the auxiliary
variable $\aux_i$. Since $\aux_i$ determines $t_i$ given $t_{<i}$, it carries the
same information as the token itself, so a model given
$\aux_i, \dots, \aux_k$ as inputs can predict $t_k$ without knowing the
intervening tokens. This is the basis of PTP: a model
$\prop(t_k \mid t_{<i}, \aux_i, \dots, \aux_k)$ emits many tokens in one call.
Figure~\ref{fig:auxiliaries} illustrates the shift of randomness from post-hoc
sampling to a random model input. We use the one-hot variant (O-PTP), whose
output probability for the proposed token doubles as a confidence estimate $c_k$.

\begin{figure}[t]
  \centering
  \includegraphics[width=0.78\textwidth]{ptp_auxiliaries.png}
  \caption{Moving the randomness. \textbf{Top:} standard decoding draws a
  distribution, then samples from it with auxiliary $\aux_i$. \textbf{Bottom:}
  PTP feeds $\aux_i$ in as an input, so the token becomes a deterministic
  function of context and auxiliaries and future tokens become jointly
  predictable. \textbf{Right:} $\aux_i$ selects a token by inverse CDF.
  Reproduced from the PTP framework~\citep{draxler2025ptp}.}
  \label{fig:auxiliaries}
\end{figure}

\subsection{What counts as a correct token}
\label{sec:correct}

Sharing the auxiliary variables gives verification a strong coupling. A proposed
token $t_k$ is \emph{correct} if it equals
$\tilde t_k = \Pick(\aux_k, \base(\cdot \mid t_{<k}))$, the token the base model
selects when the \emph{same} auxiliary is applied to its own distribution. Let
\begin{equation}
  \ncorrect = \max\{ m \geq 0 : t_{i+j} = \tilde t_{i+j} \text{ for all } j < m \}
\end{equation}
be the longest correct prefix. Because the base model evaluates
$\base(\cdot \mid t_{<k})$ at every position during verification anyway, one
token can be sampled from it for free at the first rejection, so a call commits
\begin{equation}
  \naccepted = \ncorrect + 1
  \label{eq:accepted}
\end{equation}
tokens. Under this rule the committed text is distributed exactly as the base
model's, so speedup is free of quality cost; the only question is how large
$\ncorrect$ is.

Figure~\ref{fig:partition} makes the geometry concrete for two tokens. The
auxiliary pair $(\aux_1, \aux_2)$ ranges over the unit square, and the AR
reference partitions that square into regions, one per token pair. A multi-token
model induces its own partition; where the two agree, the proposal is correct.
PTP's partition follows the reference's curved cell boundaries closely, whereas
independent per-position prediction can only produce a product partition of
axis-aligned rectangles and therefore mismatches over a large area.

\begin{figure}[t]
  \centering
  \includegraphics[width=\textwidth]{accept_partition.png}
  \caption{Verification as a partition-matching problem, for the first two
  generated tokens. The AR reference (left) partitions the auxiliary square into
  one cell per token pair. PTP (middle) reproduces the curved cell boundaries;
  independent multi-token prediction (right) can only form a product partition of
  rectangles. Correctness of a proposal is agreement of the two partitions at the
  drawn $(\aux_1, \aux_2)$. Reproduced from~\citet{draxler2025ptp}.}
  \label{fig:partition}
\end{figure}

\begin{table}[t]
  \centering
  \caption{The Spec-Bench baselines, resolved along the three axes. ``guarantee''
  is what the method promises about its output: \emph{exact} preserves the base
  model's distribution, \emph{greedy} reproduces greedy decoding only,
  \emph{biased} gives up both by design. ``overlap'' is what runs concurrently
  with verification. $p$ is the base model's distribution, $q$ / $\hat p$ the
  proposal model's, $c_j$ its confidence at node $j$, $H$ the entropy.}
  \label{tab:baselines}
  \input{tab/baselines}
\end{table}
\end{ai}

\section{The design space of multi-token decoding}
\label{sec:design}

Any multi-token decoder is characterized three main design choices: how tokens are proposed, what shape the proposal has, and which tokens get accepted into the context. The prediction model is then completed by choosing the pretrained model, the AR reference to align with, and whether any further optimizations, including call parallelization is applied.
Existing design choices are described below and existing methods are summarized in \Cref{tab:baselines}.

\subsection{Proposal Source}
\label{sec:proposal}

The framework used to sample multiple tokens in each call, builds core of each multi-token decoder. This usually involves architectural changes, training additional parameters, or finetuning the existing model. Although we will treat the proposal mechanism largely as a black-box, we will briefly describe some possible choices, summarized in \Cref{tab}

\begin{ai}

\subsection{Proposal geometry: how many chances, and how long}
\label{sec:geometry}

Given a budget of $\budget$ proposal nodes per call, the geometry question is how
to spend it: one long continuation, many short ones, or a tree. Every choice is
summarized by a \emph{width profile} $\width(d)$, the number of proposal strands
still alive at depth $d$ (\S\ref{sec:knapsack} makes this precise), and
Table~\ref{tab:geometry} classifies the options by how that profile is picked.
The column that matters is the last but one: whether the allocation accounts for
\emph{correlation} between strands. Failures are not independent --- some
positions are hard for every strand at once --- and \S\ref{sec:model} shows this
is what makes width saturate.

\begin{table}[t]
  \centering
  \caption{Proposal geometries. ``correlated'' is whether the allocation rule
  models dependence between sibling strands rather than treating per-node
  acceptance as independent. $\Sfun[w,d]$ is the survival function of
  \S\ref{sec:model}.}
  \label{tab:geometry}
  \input{tab/aspect_geometry}
\end{table}

\subsection{Verification: what counts as acceptable}
\label{sec:verification}

The exact rule of \S\ref{sec:correct} is unforgiving: a proposal is rejected
whenever the drawn auxiliary falls in the sliver where the two partitions
disagree, even if the proposed token is one the base model would have been
perfectly happy to emit at a slightly different $\aux$. On the example of
Figure~\ref{fig:accept}, exact acceptance keeps $63\%$ of the auxiliary square;
relaxing the criterion to nucleus membership keeps $99\%$ of it, with the same
proposal model and the same proposals.

\begin{figure}[t]
  \centering
  \begin{subfigure}{\textwidth}
    \includegraphics[width=\textwidth]{accept_exact.png}
    \caption{Exact acceptance: keep the proposal only where it \emph{equals} the
    reference token at the same auxiliary.}
  \end{subfigure}
  \\[0.6em]
  \begin{subfigure}{\textwidth}
    \includegraphics[width=\textwidth]{accept_nucleus.png}
    \caption{Broadened acceptance: keep the proposal wherever the reference finds
    the token plausible (nucleus membership).}
  \end{subfigure}
  \caption{The acceptance rule, not the model, decides what is kept. Both rows
  show the \emph{same} proposal partitions as Figure~\ref{fig:partition}; green
  is accepted, pink rejected. Broadening the criterion converts most of the
  rejected area into accepted area. This is the single largest lever in the
  design space, and the one that trades faithfulness for throughput most
  directly.}
  \label{fig:accept}
\end{figure}

Table~\ref{tab:verification} lists the rules. Write $q$ for the proposal model's
distribution at a position, $p$ for the base model's, $x$ for the proposed token,
$\aux$ for the shared auxiliary, and $I_p(x)$ for the inverse-CDF interval $p$
assigns to $x$. Two properties organize the column space. \emph{Does the rule
call the base model?} If not (student confidence, unconditional acceptance,
self-verification) the rule is nearly free but has no access to the distribution
we are trying to stay faithful to. \emph{Does the rule preserve that
distribution?} Only the exact and ratio rules do; the rest are deliberate
relaxations, and the empirical question is what each one costs. We implement all
eleven behind a single code path, so two runs differ only in the predicate.

\begin{table}[t]
  \centering
  \caption{Verification rules. ``base call'' is whether the rule consults the
  model being accelerated; ``output'' is the resulting guarantee. Rules with no
  entry in ``used by'' are, to our knowledge, unexplored in the multi-token
  setting; we evaluate all eleven in \S\ref{sec:experiments}.}
  \label{tab:verification}
  \input{tab/aspect_verification}
\end{table}

\subsection{Parallelization: what overlaps verification}
\label{sec:parallel}

Proposal and verification are naturally serial --- the next proposal depends on
how much of the last was accepted. Table~\ref{tab:parallel} lists the ways to
break that dependency, all of which spend compute to speculate on the
verification outcome. This axis is nearly empty in the literature: of the nine
baselines, only Lookahead overlaps anything, by fusing guessing into the same
forward pass.

\begin{table}[t]
  \centering
  \caption{Parallelization schemes. ``extra compute'' is the cost multiplier
  relative to serial drafting; $\nprop$ is the number of proposed tokens and
  $\budget$ the node budget.}
  \label{tab:parallel}
  \input{tab/aspect_parallel}
\end{table}

\subsection{What the framework generalizes}
\label{sec:generalize}

Read across the tables, the baselines are special cases.

\textbf{Geometry.} A chain is $\width(d) \equiv 1$. A static tree is a width
profile fixed offline --- Medusa picks it by greedily adding the node with the
highest calibration accuracy, Recycling hand-designs an imbalanced $81$-node,
$6$-layer shape. EAGLE-2's context-aware tree is an \emph{online} solution to
exactly the allocation problem of \S\ref{sec:knapsack}: it scores a node by
$V_i = \prod_{t_j \in \mathrm{Path}} p_j \approx \prod_j c_j$ and greedily
expands the best. That product is an independence assumption, and Medusa's
per-head calibration is the same assumption in offline form. Our model replaces
it with a correlation parameter $\rhoc$, recovers the independent case as
$\rhoc \to 0$, and solves the allocation exactly rather than greedily.

\textbf{Verification.} The exact and ratio rules are the top two rows of
Table~\ref{tab:verification}; PLD, Lookahead, Recycling and SAMD all use argmax
match. Medusa and Hydra use the entropy-adaptive threshold
$p(x) > \min(\epsilon, \delta e^{-H(p)})$, which is a genuinely different object
from a fixed threshold: it accepts liberally where the base model is uncertain
and strictly where it is confident. Our exact rule has no baseline analogue,
because it needs the auxiliary-variable coupling of \S\ref{sec:background};
speculative sampling reaches the same guarantee by a different mechanism.

\textbf{Parallelization.} Lookahead is the fused-pass row; quadratic
decoding~\citep{samragh2025future} and its budgeted variant, Partial Quadratic
Decoding~\citep{draxler2025ptp}, are the last two. \S\ref{sec:pqd} shows that the branch
weights and branch values PQD needs are both readable off the same survival
function that solves the geometry problem, which removes the two separately
calibrated estimates the original formulation requires.

\paragraph{Where the framework does not yet reach.}
Two gaps are worth stating plainly. First, our width profiles are
\emph{nonincreasing}: strands are disjoint, so a node's cost is paid at every
depth it spans. Every tree baseline instead shares prefixes, which lets
$\width(d)$ rise before it falls and changes the cost accounting to
$|s_1| + |s_2| - k$ for a shared trunk of length $k$; we return to this in
\S\ref{sec:knapsack}. Second, our budget counts nodes uniformly, whereas
retrieval-based proposals (PLD, SAMD, Recycling) cost a string match rather than
a forward pass. Both are natural extensions of the same knapsack --- a relaxed
monotonicity constraint and a per-node cost --- but we do not evaluate them here.

\section{A model of the number of correct tokens}
\label{sec:model}

The verification axis fixes \emph{when} a token is kept. The geometry axis asks
how many proposals to make and how long each should be, and answering it requires
knowing how $\ncorrect$ behaves. We therefore assume a model of $\ncorrect$ given
the number of proposal strands $K$ and the context, and we will need from it only
the survival function
\begin{equation}
  \Sfun[w, d] \;=\; \Pr\big[\, \extinct \geq d \,\big|\, w \text{ strands} \,\big],
  \qquad \extinct = \naccepted - 1,
  \label{eq:survival}
\end{equation}
where $\extinct$ is the number of generations before every strand has failed
verification and, following \eqref{eq:accepted}, $G := \naccepted = 1 + \extinct$
is the committed token count per call.

\subsection{Geometric and beta-geometric baselines}

The simplest model treats each accepted token as an independent coin flip with
success probability $p$, making $\extinct$ geometric:
\begin{equation}
  \Sfun[1, d] = p^{\,d-1},
  \qquad \Pr[\extinct = d] = (1-p)\,p^{\,d-1}.
  \label{eq:geometric}
\end{equation}
For a single strand this is already a good description:
Figure~\ref{fig:geometric} fits \eqref{eq:geometric} to $5{,}070$ observed calls
at $K=1$ with $R^2 = 0.999$. Allowing the acceptance probability to vary across
prompts, $p \sim \mathrm{Beta}(a, b)$, gives the beta-geometric model
\begin{equation}
  \Pr[\extinct = d] = \int_0^1 (1-p)\,p^{\,d-1}\,\mathrm{Beta}(p; a, b)\,\mathrm{d}p
  = \frac{B(a + d - 1,\, b + 1)}{B(a, b)},
  \label{eq:betageom}
\end{equation}
which matters because that variation is real: a likelihood-ratio test against a
single pooled $p$ rejects homogeneity decisively
($\chi^2 = 401$ on $99$ degrees of freedom, $p \approx 5 \times 10^{-38}$;
Figure~\ref{fig:heterogeneity}).

\begin{figure}[t]
  \centering
  \includegraphics[width=0.62\textwidth]{geometric.pdf}
  \caption{Committed tokens per call at $K=1$. A shifted geometric
  \eqref{eq:geometric} fitted to this width alone is already excellent; the
  branching model of \S\ref{sec:branching}, fitted \emph{jointly} across all five
  widths, gives up a little accuracy here in exchange for extrapolating across
  $K$ (Figure~\ref{fig:vsk}). Preliminary single-seed run.}
  \label{fig:geometric}
\end{figure}

\subsection{Ensembles need a correlated model}
\label{sec:branching}

Neither baseline survives the move to $K > 1$ strands. If strands failed
independently, total wipeout would vanish geometrically in $K$; empirically
$\Pr[\extinct = 1]$ falls from $41\%$ at $K = 1$ only to $13\%$ at $K = 1000$ and
flattens against a floor (Figure~\ref{fig:vsk}, left). Some positions are simply
hard, and no ensemble width rescues them.

We therefore model the strands as a branching process with a frailty that is
\emph{shared} within each generation. At depth $i$, given $\alive_{i-1} = m$
surviving strands,
\begin{equation}
  \frail_i \sim \pizero \, \delta_0 + (1 - \pizero)\,\mathrm{Beta}(\alpha, \beta),
  \qquad
  \alive_i \mid \alive_{i-1} = m,\, \frail_i \;\sim\; \mathrm{Binomial}(m, \frail_i),
  \label{eq:branching}
\end{equation}
with $\extinct = \min\{i \geq 1 : \alive_i = 0\}$. We parameterize by the
marginal per-strand success probability $p$ and the within-generation correlation
$\rhoc$ between two strands,
\begin{equation}
  \nu = \frac{1 - \rhoc}{\rhoc}, \qquad \alpha = p\,\nu, \qquad \beta = (1-p)\,\nu .
  \label{eq:reparam}
\end{equation}
Two features are load-bearing. The shared draw of $\frail_i$ correlates strands
within a generation: they tend to fail together, which is why width has
diminishing returns. The point mass $\pizero$ at $\frail = 0$ makes a fraction of
positions a guaranteed wipeout regardless of $K$, producing the observed floor;
a pure Beta-Binomial cannot. Marginalizing $\frail_i$ gives a Beta-Binomial
transition kernel
\begin{equation}
  P[m, j] = (1 - \pizero)\,\BetaBin(j \mid m, \alpha, \beta) + \pizero \, \mathbb{1}[j = 0],
  \label{eq:kernel}
\end{equation}
with $0$ absorbing, from which $\Sfun[w, \cdot]$ follows by propagation. Following
the structure of the decoder, generation $1$ is exempt from zero-inflation and
the terminal generation $\capt$ is a forced death, reflecting the hard
proposal-length cap. Fitting $(\pizero, p, \rhoc)$ by joint maximum likelihood
over all five widths simultaneously gives
$\pizero = 0.128$, $p = 0.620$, $\rhoc = 0.789$ on our checkpoint.

\begin{figure}[t]
  \centering
  \includegraphics[width=\textwidth]{hist_by_k.pdf}
  \caption{Committed tokens per call against ensemble width. Bars are observed;
  lines are the branching model \eqref{eq:branching} with a single
  $(\pizero, p, \rhoc)$ shared across all five panels (\emph{pooled}) and with
  $p$ integrated over a fitted $\mathrm{Beta}(a,b)$ population
  (\emph{heterogeneous}). $R^2$ per curve is printed in the curve's own colour.
  One three-parameter model tracks the distribution over three orders of
  magnitude of $K$.}
  \label{fig:histbyk}
\end{figure}

\begin{figure}[t]
  \centering
  \includegraphics[width=\textwidth]{vs_k.pdf}
  \caption{Why width saturates. \textbf{Left:} the probability of immediate
  wipeout decays towards an irreducible floor $\hat c = 0.118$, not to zero; a
  phenomenological $c + a K^{-b}$ captures this, and the branching model produces
  it from $\pizero$. \textbf{Right:} committed tokens per call grow roughly
  logarithmically in $K$ --- a $1000\times$ increase in proposal width buys about
  $1.5$ extra tokens.}
  \label{fig:vsk}
\end{figure}

\section{From survival to allocation: the strand-length knapsack}
\label{sec:knapsack}

Strands need not all be the same length. Let $L_s$ be the length budgeted to
strand $s$ and define the \emph{width profile}
\begin{equation}
  \width(d) := \#\{ s : L_s \geq d \},
  \label{eq:widthprofile}
\end{equation}
the number of strands still budgeted to reach depth $d$. The key structural fact
is that the shared frailty $\frail_d$ acts conditionally i.i.d.\ on every strand
present in generation $d$, so any fixed subset of a surviving population is
itself Beta-Binomial with the same $(\alpha, \beta)$ --- a thinning property.
Consequently the subpopulation that is both alive and still within its own budget
at depth $d$ is distributed exactly as a fresh homogeneous run started with
$\width(d)$ strands, \emph{independently of every other depth}, and the tail-sum
identity gives
\begin{equation}
  \mathbb{E}[\extinct] \;=\; \sum_{d \geq 1} \Pr[\extinct \geq d] \;=\; \sum_{d=1}^{\capt} \Sfun[\width(d), d].
  \label{eq:separable}
\end{equation}
Expected yield is a sum of terms each depending on the width profile at a
\emph{single} depth. That separability is what makes the design problem tractable.

Choosing strand lengths under a node budget is then
\begin{equation}
  \max_{L_1, \dots, L_K} \; \sum_{d=1}^{\capt} \Sfun[\width(d), d]
  \qquad \text{s.t.} \qquad \sum_{s} L_s \leq \budget ,
  \label{eq:knapsack}
\end{equation}
a knapsack whose items are \emph{whole strands}: buying a strand of length $L$
raises the width at every depth $1, \dots, L$ at once, so its cost is $L$ and its
value is the bundled gain across those depths.

\paragraph{Exact solution.}
The feasible width profiles are exactly the nonincreasing integer sequences
$\width(1) \geq \dots \geq \width(\capt) \geq 0$: any such sequence decomposes
into $\width(d) - \width(d+1)$ strands of length $d$, by the standard
correspondence between nonincreasing sequences and partitions, and its cost
$\sum_d \width(d)$ equals the sum of those lengths. So \eqref{eq:knapsack} is a
maximization over nonincreasing profiles, solved exactly by a dynamic program
over depth:
\begin{equation}
  V_d[w_{\mathrm{cap}}, b] \;=\; \max_{0 \leq w \leq \min(w_{\mathrm{cap}},\, b)}
  \Big\{ \Sfun[w, d] + V_{d+1}[w,\, b - w] \Big\},
  \qquad V_{\capt + 1} \equiv 0,
  \label{eq:dp}
\end{equation}
where $V_d[w_{\mathrm{cap}}, b]$ is the best achievable
$\sum_{d' \geq d} \Sfun[\width(d'), d']$ given $\width(d) \leq w_{\mathrm{cap}}$
and remaining budget $b$. Doubling the width bound until the optimum stops
touching it makes the result exact at any budget.

\begin{figure}[t]
  \centering
  \includegraphics[width=\textwidth]{width_profile.pdf}
  \caption{The allocation solved by \eqref{eq:dp}. \textbf{Left:} optimal width
  profiles at several node budgets --- wide and shallow, never uniform.
  \textbf{Right:} what the exact profile buys over spending the same budget on
  full-depth strands (the \texttt{standard} allocation). The gap is roughly
  $0.15$--$0.2$ committed tokens per call across two decades of budget.}
  \label{fig:widthprofile}
\end{figure}

\subsection{Combining with Partial Quadratic Decoding}
\label{sec:pqd}

Speculative decoding runs proposal and verification in sequence. PQD removes that
serialization by speculating on the verification outcome: while the base model
verifies $\nprop$ proposed tokens, the proposal model prepares a branch for each
possible $m \in \{0, \dots, \nprop\}$, each conditioned on ``exactly $m$ of the
proposed tokens were correct'' (Figure~\ref{fig:pqd}). When verification returns
$m^\ast$, the other branches are discarded and generation continues from the
precomputed branch $m^\ast$.

\begin{figure}[t]
  \centering
  \includegraphics[width=0.92\textwidth]{pqd_branches.png}
  \caption{Partial Quadratic Decoding. One branch is prepared per possible number
  of correct tokens, concurrently with verification; when the base model returns
  $m^\ast$, the matching branch is already extended and the rest are discarded.
  Reproduced from~\citet{draxler2025ptp}.}
  \label{fig:pqd}
\end{figure}

PQD must decide how many continuation tokens $B_m$ to give each branch under a
total budget $\budget$, which it does by maximizing an expected utility
\begin{equation}
  \max_{\sum_m B_m = \budget} \; \sum_{m} \Pr[\ncorrect = m \mid t] \; H(B_m),
  \label{eq:pqd}
\end{equation}
solved greedily by marginal gain. As originally formulated, the two ingredients
of \eqref{eq:pqd} come from different places: the branch weights are estimated by
treating errors as independent given the proposal model's confidences,
$\Pr[\ncorrect = m \mid t] \approx (1 - c_{i+m}) \prod_{k < i+m} c_k$, while the
reward $H$ is a separate monotone function calibrated on the training set.

\S\ref{sec:model} supplies both from one object. The branch weights are the
increments of the survival function,
\begin{equation}
  \Pr[\ncorrect = m] \;=\; \Sfun[K, m] - \Sfun[K, m+1],
  \label{eq:branchweights}
\end{equation}
and the value of giving a branch $n$ tokens is the expected yield of the optimal
allocation of those $n$ tokens, \ie the dynamic program \eqref{eq:dp}:
\begin{equation}
  H(n) \;=\; \max_{\width \text{ nonincreasing},\; \sum_d \width(d) \leq n} \; \sum_{d} \Sfun[\width(d), d] \;=\; V_1[\infty, n].
  \label{eq:H}
\end{equation}
Both ingredients are then consistent by construction, calibrated by the same
three parameters, and --- because $H$ is the \emph{optimal} yield rather than a
heuristic proxy --- the greedy allocation in \eqref{eq:pqd} is comparing branches
on the quantity it actually cares about. The independence approximation is
retained only where per-position confidences are genuinely more informative than
the pooled model, namely in conditioning $\Sfun$ on context; \S\ref{sec:exp:ens}
compares pooled, Bayesian-updated, and learned-head estimates of $p$ empirically.

\subsection{Worked example}
\label{sec:worked}

Table~\ref{tab:survival} gives the fitted survival function at our checkpoint's
$(\pizero, p, \rhoc) = (0.128, 0.620, 0.789)$. Two things stand out.
$\Sfun[w, 1] = 1$ for every $w \geq 1$, so width has \emph{zero} marginal value at
depth $1$ --- a single strand always reaches the first generation. And the
returns to width at a fixed depth are strongly concave: at depth $4$, going from
$1$ to $2$ strands buys more than going from $5$ to $50$.

Table~\ref{tab:knapsack} turns this into allocations. At a budget of
$\budget = 1000$ nodes --- the setting our $K = 50$ runs use --- the optimum is
$122$ strands of mixed length, from $2$ of depth $14$ down to $29$ of depth $6$,
and it yields $4.555$ committed tokens per call against $4.412$ for $50$ strands
at full depth $20$. The exact solver reallocates the deep tail of the uniform
profile into extra width at shallow depths, where the survival table says it is
worth more.

\begin{table}[t]
  \centering
  \caption{Fitted survival function $\Sfun[w, d] = \Pr[\extinct \geq d]$ for a
  homogeneous run of $w$ strands, at $(\pizero, p, \rhoc) = (0.128, 0.620,
  0.789)$. Width has no value at depth $1$ and sharply diminishing value
  thereafter.}
  \label{tab:survival}
  \input{tab/survival}
\end{table}

\begin{table}[t]
  \centering
  \caption{Optimal versus uniform allocation of a node budget, under the fitted
  model. The profile column reads $\mathrm{depth}^{\mathrm{count}}$, so
  $6^{29}$ means $29$ strands of length $6$. The exact profile trades depth for
  width and gains $\approx 0.15$ committed tokens per call.}
  \label{tab:knapsack}
  \input{tab/knapsack}
\end{table}

\section{Experiments}
\label{sec:experiments}

\subsection{Setup}

We distil a one-hot PTP student from Vicuna-7B-v1.5~\citep{zheng2023vicuna} with
gated LoRA adapters~\citep{hu2022lora} of rank $128$ and a binary auxiliary
embedding, and evaluate on the first turns of the $13$ Spec-Bench
categories~\citep{xia2024specbench} ($100$ questions), with temperature $0.7$,
top-$k$ $50$, top-$p$ $0.9$, at most $20$ proposed tokens per strand and a total
budget of $120$ proposed tokens per call, on a single NVIDIA RTX A6000.

We report three families of measure, deliberately kept separate.
\emph{Throughput per call} is $\naccepted$ averaged over calls.
\emph{Wall-clock} is milliseconds per generated token, and speedup is the AR
run's ms/token divided by the rule's --- so it charges every rule for the compute
it actually spends. \emph{Faithfulness} is measured under the base model:
$\bar q$ is the mean per-token probability the base model assigns to the
generated tokens ($\%$), PPL the corresponding perplexity, and
$p_{\mathrm{align}}$ the $p$-value of a pooled CLT $z$-test of the null
hypothesis that the generated tokens were drawn from the base distribution. A
large $p_{\mathrm{align}}$ is the operational meaning of ``indistinguishable from
the base model''.

\subsection{Acceptance rules: three orderings, not one}
\label{sec:exp:rules}

Table~\ref{tab:rules} is the main result. The headline is that the three measures
disagree.

\begin{table}[t]
  \centering
  \caption{Acceptance rules on Spec-Bench. $\bar q$ is the mean per-token
  probability under the base model, $p_{\mathrm{align}}$ the $p$-value for
  ``generated tokens came from the base distribution''. $^\dagger$
  distribution-preserving by construction. Numbers are preliminary: single-seed
  runs on $100$ Spec-Bench first turns.}
  \label{tab:rules}
  \small
  \input{tab/rules}
\end{table}

\textbf{Tokens per call is not speed.} Sequential exact PTP commits $3.63$ tokens
per call --- more than any single-call rule --- yet reaches only $1.02\times$,
because its proposal and verification passes are serialized. The single-call
variant commits $2.39$ and reaches $1.73\times$. Ensembles push tokens per call
to $4.74$ and wall-clock \emph{down} to $0.40\times$: at $K = 50$ the $1000$
proposal nodes cost more than the tokens they save. This is precisely the regime
PQD (\S\ref{sec:pqd}) exists to fix, and it is invisible if one reports only
accepted-tokens-per-call.

\textbf{Only the exact rules are faithful.} The $z$-test separates the two
classes cleanly: $p_{\mathrm{align}} = 0.25$ and $0.80$ for exact PTP, versus
$< 0.001$ for every relaxed rule. Broadening the predicate from exact matching to
nucleus $\theta = 0.98$ on the same sequential driver gains $0.67$ tokens per
call and $0.68\times$ in speedup, and costs $17.1$ points of $\bar q$
($84.7 \to 67.6$).

\textbf{Cheap rules are bad rules.} Student-confidence thresholding needs no
base-model call, and is the worst rule in the table on every axis at once: at
$\theta = 0.5$ it commits $2.01$ tokens per call, runs \emph{slower} than AR
($0.81\times$), and collapses quality to $\mathrm{PPL} = 17.0$. Accepting $k$
tokens unconditionally is the honest version of the same trade: $3.50\times$ at
$k = 4$, with $\mathrm{PPL} = 32.0$.

\textbf{An anomaly worth flagging.} The stochastic ratio rule is
distribution-preserving in theory at $\kappa = 1, \varpi = 0$, but we measure
$\bar q = 69.9$ and $p_{\mathrm{align}} < 0.001$ --- a clear deviation. We have
not marked it as lossless in Table~\ref{tab:rules} for that reason. Resolving
whether this is an artefact of residual resampling in the truncated top-$k$ space
or a genuine implementation defect is left to future work; we report it rather
than assume the theory.

\begin{figure}[t]
  \centering
  \includegraphics[width=\textwidth]{pareto.pdf}
  \caption{Throughput against faithfulness for all $62$ configurations. The two
  panels share a $y$-axis but order the rules differently: high tokens per call
  (left) does not imply wall-clock speedup (right). Only the exact rules sit on
  the autoregressive reference line.}
  \label{fig:pareto}
\end{figure}

\subsection{Threshold sweeps}

Figure~\ref{fig:sweeps} sweeps $\theta$ for the three threshold-style families.
Base-model-probability thresholding degrades gracefully and monotonically: as
$\theta \to 0$ it accepts everything, reaching $17$ tokens per call and
$4.1\times$, and as $\theta$ grows it recovers the base model's quality. Nucleus
acceptance occupies a narrow, useful band near $\theta \approx 0.9$--$0.98$.
Student-confidence thresholding is non-monotone and dominated everywhere.

\begin{figure}[t]
  \centering
  \includegraphics[width=\textwidth]{sweeps.pdf}
  \caption{Threshold sweeps. Separate panels rather than twin axes, because the
  three measures are on unrelated scales. Dashed lines mark the autoregressive
  reference.}
  \label{fig:sweeps}
\end{figure}

\subsection{Ensembles and allocation}
\label{sec:exp:ens}

Table~\ref{tab:ensemble} varies how the per-strand acceptance probability driving
the allocation of \S\ref{sec:knapsack} is estimated: a single pooled $p$, an
online Bayesian update (conjugate Beta-Binomial, or a grid reweighted by the
survival likelihood), and a small learned head predicting $p$ from the hidden
state. Differences are modest at $K = 5$ and grow at $K = 50$, consistent with
the allocation mattering more when there is more budget to misallocate.

\begin{table}[t]
  \centering
  \caption{Estimating the allocation's acceptance probability. All rows use the
  nucleus criterion at the stated $\theta$; ``allocation'' is how $p$ is
  estimated for the knapsack of \S\ref{sec:knapsack}. Preliminary.}
  \label{tab:ensemble}
  \small
  \input{tab/ensemble}
\end{table}

\subsection{Does the model predict the data?}

Table~\ref{tab:fit} and Figure~\ref{fig:histbyk} validate \S\ref{sec:model}. One
three-parameter model, fitted jointly, reproduces the per-call histogram at every
width with $R^2 \in [0.976, 0.996]$ and predicts $\mathbb{E}[G]$ to within
$0.08$ tokens at $K \leq 50$. It degrades at $K = 1000$ (predicting $5.09$ against
an observed $4.81$), which is the extrapolation limit of a depth-uniform
$(\pizero, p, \rhoc)$. Integrating $p$ over the fitted population
\eqref{eq:betageom} improves the fit at every width, by between $+0.3$ and
$+2.0$ points of $R^2$, but not monotonically in $K$: the largest gain is at
$K = 50$ and the gain at $K = 1000$ is smaller than at $K = 1$
(Figure~\ref{fig:heterogeneity}). Prompt-level heterogeneity is therefore real
and worth modelling, but it is not the dominant source of the pooled model's
residual error, and the $K = 1000$ shortfall is not explained by it.

\begin{table}[t]
  \centering
  \caption{Model validation across ensemble widths. $R^2$ of the predicted
  per-call histogram for the pooled and heterogeneous models, and the
  likelihood-ratio statistic against a single pooled $p$ (99 d.o.f.).}
  \label{tab:fit}
  \small
  \input{tab/fit}
\end{table}

\begin{figure}[t]
  \centering
  \includegraphics[width=\textwidth]{heterogeneity.pdf}
  \caption{Prompt-level heterogeneity. \textbf{Left:} per-question acceptance
  probabilities at $K=1$, the fitted population $\mathrm{Beta}(a,b)$, and the
  pooled estimate; the empirical spread is wider than the population because the
  per-question MLEs carry estimation noise. \textbf{Right:} the resulting gain in
  histogram fit, recomputed against the same predictions used in
  Figure~\ref{fig:histbyk} and Table~\ref{tab:fit}. The gain is positive at
  every width but not monotone in $K$.}
  \label{fig:heterogeneity}
\end{figure}

\subsection{Self-verification}

Dropping the base model entirely and verifying against the proposal model's own
autoregressive prediction is attractive --- one set of weights, no second model
--- and it does raise tokens per call ($3.86$ exact, $6.45$ with nucleus
$\theta = 0.8$). But faithfulness collapses: $\bar q = 65.9$ for exact
self-verification and $24.7$ with nucleus. This is expected rather than
surprising. Self-verification removes the only reference to the distribution we
are asking the output to match, so the generation is free to drift, and the
metric that detects drift duly does. Self-verification is the right tool when the
goal is a fast standalone generator, and the wrong one when the goal is to
accelerate a \emph{specific} base model.

\section{Related work}

Table~\ref{tab:baselines} already resolves the methods we compare against along
our three axes, so we note here only what they have in common. Speculative
decoding established the proposal-and-verify
template~\citep{leviathan2023speculative,chen2023accelerating}, and the work
since has almost entirely gone into the \emph{proposal}: extra decoding
heads~\citep{stern2018blockwise,cai2024medusa}, sequentially dependent
heads~\citep{ankner2024hydra}, feature-space drafting with static and then
dynamic draft trees~\citep{li2024eagle,li2024eagle2}, Jacobi-style parallel
decoding~\citep{fu2024lookahead}, and retrieval from the
prompt~\citep{saxena2023pld}. Multi-token prediction has also been used as a
training objective in its own right~\citep{gloeckle2024multitoken}. Verification
is correspondingly under-explored: the acceptance rule is inherited, not chosen,
with Medusa's entropy-adaptive threshold the notable exception. And the standard
benchmark protocol reports accepted-tokens-per-call and
speedup~\citep{xia2024specbench} without any test of distributional
faithfulness, which \S\ref{sec:exp:rules} shows is where the rules actually
differ. Allocation under a draft budget is treated heuristically throughout ---
EAGLE-2's greedy expansion is the most explicit version, and it assumes
independence between sibling nodes. PTP~\citep{draxler2025ptp} supplies the
proposal mechanism we build on, and Partial Quadratic Decoding, which budgets the
quadratic expansion of~\citet{samragh2025future}; the connection
between auxiliary variables and inverse autoregressive
flows~\citep{kingma2016iaf} is noted there. Our branching model is a
zero-inflated, shared-frailty Galton--Watson process~\citep{athreya1972branching}.

\section{Conclusion}

Treating the acceptance rule as a design variable rather than a fixed detail
changes the picture of multi-token inference. Accepted-tokens-per-call,
wall-clock speedup, and faithfulness to the base model rank decoding rules
differently, and reporting only the first --- as is common --- can recommend a
configuration that is both slower and less faithful than the alternative. A
three-parameter branching model of the number of correct tokens per call captures
the per-call distribution across three orders of magnitude of proposal width, and
makes the proposal-geometry question an exactly solvable knapsack whose solution
also supplies the branch weights and branch values that Partial Quadratic
Decoding needs.

\paragraph{Limitations.}
All numbers are preliminary, single-seed, and from one model pair
(Vicuna-7B and one distilled student) on one GPU; the wall-clock conclusions in
particular are hardware- and implementation-specific, and the ensemble rules
would be read differently on hardware with spare parallel capacity. We evaluate
only first turns of Spec-Bench, which leaves no reference completion, so our
quality measures are all relative to the base model rather than to a ground
truth; a rule that drifts towards \emph{better} text than the base model's would
be penalized by $p_{\mathrm{align}}$ exactly as one that drifts towards worse. The
fitted $(\pizero, p, \rhoc)$ is depth-uniform and checkpoint-specific, and its
extrapolation degrades at $K = 1000$. The ratio-rule anomaly of
\S\ref{sec:exp:rules} is unresolved.
\end{ai}

\bibliography{iclr2026_conference}
\bibliographystyle{iclr2026_conference}

\appendix
\section{All decoding configurations}

Table~\ref{tab:all} lists every configuration, sorted by wall-clock speedup.
Re-run variants of identical configurations are excluded, as are runs marked
superseded in the source tree. ``gzip'' is the compression ratio of the generated
text, a model-free repetition proxy; lower means more repetitive.

\begin{table}[h]
  \centering
  \caption{All $62$ decoding configurations, sorted by speedup. Preliminary.}
  \label{tab:all}
  \tiny
  \input{tab/all_runs}
\end{table}

\end{document}
